% Vietnamese translation for OI Public Library
% Source-PDF: IOI中国国家候选队论文/2004/周源.pdf
\documentclass[11pt]{article}
\usepackage{oipl-vn}
\usepackage{tikz}
\usetikzlibrary{arrows.meta,calc,patterns}

\title{Bàn về ứng dụng tư tưởng kết hợp số và hình trong thi tin học}
\author{Zhou Yuan}
\date{}

\begin{document}
\maketitle

\origtitle{On the application of numeric-geometric thinking in informatics contests}
\sourcepath{IOI China candidate team papers / 2004 / Zhou Yuan.pdf}

\section*{Tóm tắt}

Số và hình là hai đối tượng cổ xưa nhất, cơ bản nhất của toán học; tư tưởng
kết hợp số và hình cũng là một tư tưởng toán học quan trọng.

Bài viết này lấy một số bài toán trong Olympic Tin học làm ví dụ, rồi bàn về
triển vọng ứng dụng rộng rãi của tư tưởng kết hợp số và hình trong giải bài,
theo hai hướng nhấn mạnh khác nhau: dùng hình hỗ trợ số và dùng số hỗ trợ
hình.

Phần cuối phân tích hai đặc tính quan trọng của tư tưởng này, từ đó chỉ ra
rằng trọng tâm của ``kết hợp số và hình'' nằm ở chữ kết hợp một cách hữu cơ.
Hy vọng bài viết giúp người đọc vận dụng linh hoạt tư tưởng ấy trong thi đấu
thực tế.

\paragraph{Từ khóa.}
Thi tin học; tư tưởng toán học; kết hợp số và hình; dùng số hỗ trợ hình; dùng
hình hỗ trợ số; mâu thuẫn biện chứng; tính đa dạng; khác biệt cá nhân; tư
duy; lập trình; độ phức tạp thời gian và không gian.

\tableofcontents

\section{Dẫn nhập}

Số và hình là hai đối tượng cổ xưa nhất, cơ bản nhất của toán học; tư tưởng
kết hợp số và hình cũng là một tư tưởng toán học quan trọng.

Trong các kỳ thi tin học hiện nay, đằng sau một số đề bài rối rắm thường ẩn
chứa bối cảnh hình học phong phú; ngược lại, nhiều bài toán hình học tính
toán lại cần dựa vào năng lực tính toán số thực mạnh mẽ của máy tính. Như
nhận xét nổi tiếng của Hoa La Canh rằng kết hợp số và hình đem lại vô vàn
lợi ích, trong thiết kế thuật toán và chương trình, nếu vận dụng khéo léo tư
tưởng này, ta có thể phá được nút thắt của bài toán, biến khó thành dễ và tìm
ra đường vào cho lời giải.

Tư tưởng kết hợp số và hình thường bao gồm hai mặt: dùng hình hỗ trợ số và
dùng số hỗ trợ hình.

\section{Dùng hình hỗ trợ số}

Như đã nói ở trên, phía sau những quan hệ đại số phức tạp trong một số đề bài
thường có một bối cảnh hình học giàu cấu trúc. Nhờ tính chất của hình vẽ,
những quan hệ số lượng và khái niệm trừu tượng vốn khó nắm bắt có thể trở nên
trực quan hơn, qua đó mở ra lối tắt để thiết kế thuật toán.

\subsection{Ví dụ 1: chứng minh bổ đề Raney}

\subsubsection*{Phát biểu bài toán}

Cho dãy số nguyên
\[
  A=\{A_i\mid i=1,2,\ldots,N\},
\]
và các tổng tiền tố $S_k=A_1+\cdots+A_k$. Giả sử tổng của toàn bộ dãy là
$S_N=1$.

Chứng minh rằng trong $N$ biểu diễn vòng của $A$ có đúng một dãy $B$ sao cho
mọi tổng tiền tố của $B$ đều dương. Ở đây, có thể tưởng tượng dãy ban đầu nằm
trên một vòng tròn; cắt vòng tại một vị trí bất kỳ sẽ thu được một biểu diễn
vòng tuyến tính.

\subsubsection*{Phân tích}

Xét trước một ví dụ:
\[
  A=\langle 1,4,-5,3,-2,0\rangle .
\]
Sáu biểu diễn vòng của nó là
\[
\begin{array}{rcl}
1&:&\langle 1,4,-5,3,-2,0\rangle,\\
2&:&\langle 4,-5,3,-2,0,1\rangle,\\
3&:&\langle -5,3,-2,0,1,4\rangle,\\
4&:&\langle 3,-2,0,1,4,-5\rangle,\\
5&:&\langle -2,0,1,4,-5,3\rangle,\\
6&:&\langle 0,1,4,-5,3,-2\rangle.
\end{array}
\]
Trong đó chỉ dãy thứ tư có các tổng tiền tố
$3,1,1,2,6,1$, tất cả đều dương.

Nếu dùng phương pháp đại số hoặc tổ hợp thông thường để chứng minh kết luận
thú vị này, ta khó thấy nên bắt đầu từ đâu. Nhưng nếu để ``hình'' tham gia
vào suy nghĩ, vấn đề trở nên rất đơn giản.

\paragraph{Đưa mục tiêu về hình học.}
Mở rộng dãy $A$ theo chu kỳ thành dãy vô hạn
\[
  \langle A_1,A_2,\ldots,A_N,A_1,A_2,\ldots,A_N,\ldots\rangle
\]
và đồng thời xét các tổng tiền tố $S_i$ của dãy vô hạn đó. Vì dãy có chu kỳ
và tổng mỗi chu kỳ bằng $1$, với mọi $k>0$ ta có
\[
  S_{k+N}=S_k+1.
\]
Nếu vẽ đồ thị của hàm tổng tiền tố, ta có thể nói hàm này có ``độ dốc trung
bình'' bằng $1/N$: đi sang phải $N$ đơn vị thì giá trị hàm tăng thêm $1$.
Do đó có thể kẹp toàn bộ các điểm của đồ thị giữa hai đường thẳng song song
có hệ số góc $1/N$.

\begin{center}
\begin{tikzpicture}[x=0.67cm,y=0.55cm, every node/.style={font=\small}]
  \draw[-{Stealth[length=2mm]}] (-0.2,0) -- (20.2,0) node[right] {$k$};
  \draw[-{Stealth[length=2mm]}] (0,-3.2) -- (0,6.7) node[above] {$S_k$};
  \foreach \x in {0,6,12,18} \draw[gray!35] (\x,-3.1) -- (\x,6.2);
  \foreach \y in {-2,0,2,4,6} \draw[gray!20] (0,\y) -- (19.5,\y);

  \coordinate (p0) at (0,0);
  \coordinate (p1) at (1,1);
  \coordinate (p2) at (2,5);
  \coordinate (p3) at (3,0);
  \coordinate (p4) at (4,3);
  \coordinate (p5) at (5,1);
  \coordinate (p6) at (6,1);
  \coordinate (p7) at (7,2);
  \coordinate (p8) at (8,6);
  \coordinate (p9) at (9,1);
  \coordinate (p10) at (10,4);
  \coordinate (p11) at (11,2);
  \coordinate (p12) at (12,2);
  \coordinate (p13) at (13,3);
  \coordinate (p14) at (14,7);
  \coordinate (p15) at (15,2);
  \coordinate (p16) at (16,5);
  \coordinate (p17) at (17,3);
  \coordinate (p18) at (18,3);
  \draw[thick,blue!60!black] (p0)--(p1)--(p2)--(p3)--(p4)--(p5)--(p6)--(p7)--(p8)--(p9)--(p10)--(p11)--(p12)--(p13)--(p14)--(p15)--(p16)--(p17)--(p18);
  \foreach \i in {0,...,18} \fill[blue!60!black] (p\i) circle (1.3pt);

  \draw[red!65!black,thick] (0,-0.5) -- (19,2.666);
  \draw[red!65!black,thick,dashed] (0,4.5) -- (12,6.5);
  \node[red!65!black] at (15.2,2.05) {hệ số góc $1/N$};
  \fill[red!65!black] (3,0) circle (2.3pt);
  \fill[red!65!black] (9,1) circle (2.3pt);
  \fill[red!65!black] (15,2) circle (2.3pt);
  \node[below] at (3,-0.2) {$m$};
  \node[below] at (9,-0.2) {$m+N$};
  \node[below] at (15,-0.2) {$m+2N$};
  \node at (9,-2.55) {Đồ thị tổng tiền tố của dãy mở rộng chu kỳ, minh họa với $N=6$.};
\end{tikzpicture}
\end{center}

Chú ý đường thẳng thấp hơn. Trong mỗi đoạn dài $N$ liên tiếp theo trục hoành,
nó cắt đồ thị tổng tiền tố đúng một lần, vì một đường thẳng có hệ số góc
$1/N$ nhiều nhất đi qua một điểm nguyên trong mỗi đoạn như vậy. Giao điểm đó
là điểm thấp nhất trong $N$ điểm tiếp theo. Vì thế nếu cắt dãy ngay sau vị
trí ấy, mọi tổng tiền tố thu được đều dương. Đồng thời, việc mỗi đoạn dài
$N$ chỉ có một giao điểm cũng chứng minh tính duy nhất của biểu diễn vòng cần
tìm.

\paragraph{Tiểu kết.}
Một lập luận hình học rất ngắn đã chứng minh được bổ đề Raney nổi tiếng; sự
gọn gàng này khó đạt được bằng các cách khác.

Bổ đề Raney có nhiều ứng dụng. Các công thức tổ hợp của số Catalan và các mở
rộng của số Catalan đều có thể xử lý nhẹ nhàng bằng bổ đề này. Chẳng hạn bài
\textit{Sequence} trong ngày thi thứ hai của kỳ tuyển chọn Thượng Hải năm
2003, cũng như bài \textit{Necklace} trong một vòng luyện tập OIBH, đều có
thể giải rất tự nhiên bằng bổ đề Raney. Phần phụ lục trình bày lời giải cho
\textit{Sequence}.

Dùng hình học để hỗ trợ tư duy không chỉ xuất hiện trong đếm tổ hợp. Nó còn
thấm vào nhiều nhánh của thiết kế và tối ưu thuật toán. Phương pháp ``tối ưu
theo hệ số góc'', vốn phổ biến trong những năm gần đây, là một ví dụ rất tốt
khác.

\subsection{Ví dụ 2: bài toán giá trị trung bình lớn nhất}

Bài toán này xuất hiện trong USACO March Open 2003.

\subsubsection*{Phát biểu bài toán}

Cho một dãy số dương $a_1,a_2,\ldots,a_N$ và một số $F$. Định nghĩa
\[
  \operatorname{ave}(i,j)=
  \frac{a_i+\cdots+a_j}{j-i+1},\qquad i\le j.
\]
Hãy tìm
\[
  \max\{\operatorname{ave}(a,b)\mid 1\le a\le b\le N,\ b-a+1\ge F\},
\]
tức là tìm đoạn con có độ dài ít nhất $F$ và có trung bình cộng lớn nhất.
Giới hạn: $F\le N\le 10^5$.

\subsubsection*{Phân tích}

Một thuật toán liệt kê đơn giản sẽ duyệt mọi cặp $(a,b)$ thỏa điều kiện, tính
\(\operatorname{ave}(a,b)\), rồi cập nhật đáp án hiện tại. Nhưng $N$ rất lớn,
nên thuật toán bậc hai hiển nhiên không dùng được. Câu hỏi là liệu ta có thể
tối ưu cách liệt kê kém hiệu quả đó hay không. Câu trả lời là có.

\paragraph{Đưa mục tiêu về hình học.}
Trước hết đặt tổng tiền tố
\[
  S_i=a_1+a_2+\cdots+a_i,\qquad S_0=0.
\]
Khi đó
\[
  \operatorname{ave}(i,j)=\frac{S_j-S_{i-1}}{j-(i-1)}.
\]
Nếu vẽ các điểm
\[
  P_i=(i,S_i),\qquad 0\le i\le N,
\]
trên mặt phẳng tọa độ, biểu thức trên chính là hệ số góc của đường thẳng đi
qua $P_j$ và $P_{i-1}$. Vì vậy bài toán trở thành: cho $N+1$ điểm trên mặt
phẳng, tìm đường nối hai điểm có khoảng cách ngang ít nhất $F$ và có hệ số
góc lớn nhất.

Ta quy ước rằng với $i<j$, chỉ xét đường nối từ điểm sau $P_j$ về điểm trước
$P_i$; không xét chiều ngược lại. Nói cách khác, với mỗi điểm chỉ kiểm tra
các điểm nằm đủ xa phía trước nó. Định nghĩa tập kiểm tra của $P_i$ là
\[
  G_i=\{P_j\mid 0\le j\le i-F\}.
\]
Khi $i<F$ thì $G_i$ rỗng. Ý nghĩa rất rõ: trong thuật toán bậc hai, nếu cần
kiểm tra $\operatorname{ave}(a,b)$ thì chắc chắn $P_{a-1}\in G_b$. Vì thế
thuật toán bậc hai cũng có thể mô tả là lần lượt liệt kê điểm đang kiểm tra
$P_b$, sau đó liệt kê từng điểm bị kiểm tra $P_a\in G_b$ và xét hệ số góc
$k(P_aP_b)$.

\paragraph{Xây dựng đường gấp khúc lồi dưới.}
Giả sử trong một tập kiểm tra có ba điểm $P_i,P_j,P_k$ với $i<j<k$, và
$P_j$ nằm trên phần lồi phía trên của đoạn $P_iP_k$, tức là
\[
  k(P_i,P_j)>k(P_j,P_k).
\]
Khi đó có thể chứng minh $P_j$ là điểm thừa.

\begin{center}
\begin{tikzpicture}[x=0.82cm,y=0.72cm, every node/.style={font=\small}]
  \draw[-{Stealth[length=2mm]}] (-0.4,0) -- (10.2,0) node[right] {$x$};
  \draw[-{Stealth[length=2mm]}] (0,-0.3) -- (0,6.5) node[above] {$y$};
  \coordinate (Pi) at (2.2,1.4);
  \coordinate (Pj) at (5.0,3.0);
  \coordinate (Pk) at (7.7,4.75);
  \draw[thick] (Pi)--(Pj)--(Pk);
  \draw[dashed] (Pi)--(Pk);
  \draw[fill=gray!18,draw=none] (Pi)--(Pj)--(5.0,6.2)--(2.2,6.2)--cycle;
  \draw[fill=gray!30,draw=none] (Pj)--(Pk)--(7.7,0.3)--(5.0,0.3)--cycle;
  \draw[thick] (Pi)--(Pj)--(Pk);
  \fill (Pi) circle (2pt) node[below left] {$P_i$};
  \fill (Pj) circle (2pt) node[above left] {$P_j$};
  \fill (Pk) circle (2pt) node[above right] {$P_k$};
  \node at (3.1,5.35) {vùng 1};
  \node at (6.4,1.0) {vùng 2};
  \node at (4.75,2.0) {phần giao};
  \node at (5,-0.85) {Nếu $P_tP_j$ thắng cả $P_tP_i$ và $P_tP_k$, $P_t$ phải rơi vào phần giao bất khả thi.};
\end{tikzpicture}
\end{center}

Thật vậy, nếu $k(P_t,P_j)>k(P_t,P_i)$ thì điểm $P_t$ phải nằm phía trên
đường $P_iP_j$, tức trong vùng 1. Tương tự, nếu
$k(P_t,P_j)>k(P_t,P_k)$ thì $P_t$ phải nằm phía dưới đường $P_jP_k$, tức
trong vùng 2. Nếu hệ số góc của $P_tP_j$ đồng thời lớn hơn hệ số góc của
$P_tP_i$ và $P_tP_k$, điểm $P_t$ phải nằm trong phần giao của hai vùng ấy.
Nhưng phần giao này không thỏa giả thiết ban đầu $t>j$. Do đó khi $P_t$ nằm
ở bất kỳ vị trí hợp lệ nào, hệ số góc $P_tP_j$ luôn không thể là lớn nhất:
nó sẽ thua một trong hai đường qua $P_i$ hoặc $P_k$. Điểm $P_j$ vì vậy có
thể loại khỏi tập kiểm tra mà không ảnh hưởng đáp án.

Kết luận trên cho biết: trong tập kiểm tra của mọi điểm $P_t$, một điểm lồi
phía trên sẽ không đóng góp cho nghiệm tối ưu. Nếu loại mọi điểm như vậy, các
điểm còn lại tạo thành một đường gấp khúc lồi dưới. Bây giờ ta chỉ cần tìm
trên đường gấp khúc lồi dưới ấy một điểm sao cho đường nối với $P_t$ có hệ
số góc lớn nhất. Rõ ràng đường đó đạt cực đại khi nó tiếp xúc với đường gấp
khúc.

\begin{center}
\begin{tikzpicture}[x=0.82cm,y=0.75cm, every node/.style={font=\small}]
  \draw[-{Stealth[length=2mm]}] (-0.3,0) -- (10.2,0) node[right] {$x$};
  \draw[-{Stealth[length=2mm]}] (0,-0.3) -- (0,6.5) node[above] {$y$};
  \coordinate (A) at (1.2,3.4);
  \coordinate (B) at (2.6,1.7);
  \coordinate (C) at (4.6,1.25);
  \coordinate (D) at (6.5,1.8);
  \coordinate (E) at (8.0,3.1);
  \coordinate (T) at (8.9,5.8);
  \draw[thick,blue!60!black] (A)--(B)--(C)--(D)--(E);
  \foreach \p/\name in {A/{},B/{},C/{Q},D/{},E/{}} {
    \fill[blue!60!black] (\p) circle (1.8pt);
  }
  \fill[red!70!black] (T) circle (2pt) node[above] {$P_t$};
  \draw[red!70!black,thick] (T)--(C);
  \draw[gray!60,dashed] (T)--(B);
  \draw[gray!60,dashed] (T)--(D);
  \node[below] at (C) {$Q$};
  \node at (5.2,-0.85) {Đường qua $P_t$ có hệ số góc lớn nhất tại điểm tiếp xúc $Q$.};
\end{tikzpicture}
\end{center}

\paragraph{Duy trì đường gấp khúc lồi dưới.}
Mục tiêu của phần này là lấy được đường gấp khúc lồi dưới cho từng điểm kiểm
tra với chi phí nhỏ nhất có thể.

Thuật toán bắt đầu từ $P_F$. Đây là điểm ngoài cùng bên trái có tập kiểm tra
không rỗng; tập ấy chỉ chứa $P_0$, hiển nhiên thỏa điều kiện lồi dưới. Sau
đó ta lần lượt xét các điểm mới $P_{F+1},P_{F+2},\ldots,P_N$.

Trong quá trình xét, ta duy trì đường gấp khúc lồi dưới. Mỗi khi đến điểm
mới $P_t$, ta thêm điểm $P_{t-F}$ vào đầu phải của đường gấp khúc. Điểm mới
có thể làm một vài điểm ở đầu phải trở thành điểm lồi phía trên. Dùng thủ
tục tương tự dựng bao lồi, ta xóa dần các điểm này để bảo đảm tính lồi dưới.
Vì mỗi điểm chỉ được thêm một lần và bị xóa nhiều nhất một lần, chi phí duy
trì bình quân cho mỗi bước là $O(1)$. Như vậy trong tổng thời gian $O(N)$ ta
thu được đường gấp khúc lồi dưới của mọi tập kiểm tra.

Một mô tả ngắn gọn của phần duy trì là:
\[
\begin{array}{l}
\text{với }t=F,F+1,\ldots,N\text{:}\\
\quad \text{chèn }P_{t-F}\text{ vào cuối hàng đợi lồi;}\\
\quad \text{trong khi ba điểm cuối không còn tạo lồi dưới, xóa điểm giữa;}\\
\quad \text{dùng hàng đợi này để tìm điểm cho hệ số góc lớn nhất với }P_t.
\end{array}
\]

\paragraph{Tối ưu cuối cùng: dùng tính đơn điệu của hình vẽ.}
Vấn đề còn lại là tìm đường qua $P_t$ tiếp xúc với đường gấp khúc. Một cách
trực tiếp là tìm nhị phân, mất $O(\log N)$ cho mỗi điểm. Nhưng từ tính chất
hình học ta có cách đơn giản và nhanh hơn. Với mỗi điểm trên đường gấp khúc,
phạm vi hệ số góc của các tiếp tuyến qua nó đã được xác định; hơn nữa theo
tính đơn điệu của hệ số góc trên hàm lồi dưới, nếu khi xét $P_t$ ta đã tìm
được một điểm tiếp xúc $A$ trên đường gấp khúc, thì mọi điểm trước $A$ đều có
thể bỏ khỏi hàng đợi. Các tiếp tuyến qua những điểm ấy có hệ số góc nhỏ hơn
nghiệm tốt nhất đã biết và sẽ không tạo đóng góp lớn hơn.

Vì thế ta giữ thêm một con trỏ chỉ tiến về sau để tìm tiếp điểm. Trong cài
đặt, khi điểm kế tiếp trên đường gấp khúc cho hệ số góc với $P_t$ lớn hơn
hoặc bằng điểm hiện tại, ta tăng con trỏ. Con trỏ không bao giờ lùi, nên chi
phí bình quân cũng là $O(1)$.

Đến đây, bài toán được giải trọn vẹn với độ phức tạp thời gian và không gian
đều là $O(N)$.

\paragraph{Tiểu kết.}
Nhìn lại quá trình giải, ngay từ đầu ta đã chọn mặt phẳng hình học làm công
cụ tư duy. Nhờ đó rất nhanh chóng phát hiện kết luận quan trọng: những điểm
có thể đóng góp cho nghiệm tối ưu trong tập kiểm tra tạo thành một hàm lồi
dưới. Sau đó ta dùng phương pháp dựng bao lồi của hình học tính toán để duy
trì đường gấp khúc lồi dưới, rồi tiếp tục dựa vào tính đơn điệu của hệ số góc
trên hàm lồi dưới để tìm tiếp tuyến bằng con trỏ. Toàn bộ lời giải xoay quanh
trục hình học phẳng, với hệ số góc làm sợi chỉ chính, tránh được các biến đổi
đại số dễ gây rối mắt. Đây là một ví dụ kinh điển của việc dùng hình hỗ trợ
số.

Nhân tiện, phương pháp kiểu này còn có thể dùng để tăng tốc quá trình ra
quyết định trong nhiều thuật toán quy hoạch động. Các bài như \textit{Batch}
của IOI 2002 hay \textit{Euro} của BOI 2003 đều có thể vận dụng tư tưởng
``tối ưu hệ số góc'' để cải thiện hiệu quả. Những đặc điểm chung của lớp bài
này rất đáng nghiên cứu thêm, nhưng không thuộc phạm vi bài viết.

\section{Dùng số hỗ trợ hình}

Pythagoras của Hy Lạp cổ đại cho rằng ``vạn vật đều là số''. Quả thật, số là
một trong những cách tốt nhất để phản ánh đặc trưng bản chất của sự vật.
Trong lịch sử phát triển toán học, chính sau khi hình học giải tích ra đời,
con người mới hiểu sâu hơn về nhiều loại đường cong phức tạp. Trong thời đại
thông tin ngày nay, máy tính xử lý mọi loại đối tượng rốt cuộc cũng quy về
các phép toán cơ bản trên số nhị phân; từ đó có thể thấy tầm quan trọng của
số.

Trong thi tin học hiện nay, một số đề bài mô tả hình vẽ rất phức tạp, dễ làm
thí sinh lạc vào mê trận. Khi ấy, dùng số hỗ trợ hình, tức lập tức nắm lấy
đặc trưng bản chất của hình bằng biểu diễn số, thường là một phương pháp giải
tốt.

\subsection{Ví dụ 3: xưởng vẽ}

Bài toán này là \textit{Painter's Studio} của POI, vòng I, lần V.

\subsubsection*{Phát biểu bài toán}

Định nghĩa ma trận vuông kích thước $0$ là một ma trận $1\times 1$, trong ô
duy nhất có một lỗ nhỏ. Với $i>0$, ma trận vuông kích thước $i$ được định
nghĩa đệ quy như sau: chia hình vuông thành bốn phần bằng nhau; ô phần tư
trên trái không có lỗ, còn ba phần tư còn lại là ba bản sao của ma trận
kích thước $i-1$.

\begin{center}
\begin{tikzpicture}[x=0.75cm,y=0.75cm, every node/.style={font=\small}]
  \foreach \x/\y in {0/0,1/0,0/1,1/1} {
    \draw[gray!45] (\x,\y) rectangle ++(1,1);
  }
  \fill[gray!20] (0,1) rectangle ++(1,1);
  \node at (0.5,1.5) {trống};
  \node at (1.5,1.5) {$M_{i-1}$};
  \node at (0.5,0.5) {$M_{i-1}$};
  \node at (1.5,0.5) {$M_{i-1}$};
  \draw[thick] (0,0) rectangle (2,2);
  \node at (1,-0.55) {Cấu trúc đệ quy của $M_i$.};

  \begin{scope}[xshift=6.2cm]
    \draw[step=0.5,gray!35] (0,0) grid (4,4);
    \draw[thick] (0,0) rectangle (4,4);
    \foreach \x/\y in {
      0/0,1/0,2/0,3/0,
      0/1,2/1,3/1,
      2/2,3/2,
      0/2,1/2,3/3
    } \fill[blue!60!black] (\x+0.25,\y+0.25) circle (1.7pt);
    \draw[red!65!black,thick] (1.0,1.0) rectangle (5.0,5.0);
    \foreach \x/\y in {
      1/1,2/1,3/1,4/1,
      1/2,3/2,4/2,
      3/3,4/3,
      1/3,2/3,4/4
    } \fill[red!65!black] (\x+0.25,\y+0.25) circle (1.7pt);
    \node at (2,-0.55) {Hai bản $N=2$ lệch nhau $(X,Y)=(2,2)$ có 3 lỗ chung.};
  \end{scope}
\end{tikzpicture}
\end{center}

Cho kích thước $N$ của ma trận và hai tham số $X,Y$. Lấy hai ma trận cùng
kích thước $N$, đặt một ma trận chồng lên ma trận còn lại, rồi dời ma trận
phía trên sang phải $X$ cột và lên trên $Y$ hàng. Hỏi sau thao tác này hai
ma trận có bao nhiêu lỗ nhỏ chung. Giới hạn: $N\le 100$.

\subsubsection*{Phân tích}

Phân tích trực tiếp hình giao của hai ma trận là có thể làm được. Tác giả
từng xem một số báo cáo lời giải của các anh chị đội tuyển, và phần lớn đi
theo hướng ấy. Nhưng cách làm đó rất rườm rà, đòi hỏi lượng suy nghĩ lớn:
phải đánh dấu vùng giao nhau của hai hình rồi xét nhiều trường hợp. Có thể
thấy nếu đối đầu trực tiếp với ``hình'', con đường lời giải sẽ khá gập ghềnh.

\paragraph{Đưa mục tiêu về số.}
Ta chuyển sang mặt còn lại, tức ``số''. Đánh số hàng và cột từ $0$ đến
$2^N-1$ theo hai chiều $x,y$. Khi đó mỗi ô có một tọa độ duy nhất $(x,y)$.

\begin{center}
\begin{tikzpicture}[x=0.55cm,y=0.55cm, every node/.style={font=\small}]
  \draw[step=1,gray!35] (0,0) grid (8,8);
  \draw[thick] (0,0) rectangle (8,8);
  \draw[-{Stealth[length=2mm]}] (0,-0.6) -- (8.6,-0.6) node[right] {$x$};
  \draw[-{Stealth[length=2mm]}] (-0.6,0) -- (-0.6,8.6) node[above] {$y$};
  \foreach \x in {0,1,...,7} \node[below] at (\x+0.5,-0.6) {\scriptsize \x};
  \foreach \y in {0,1,...,7} \node[left] at (-0.6,\y+0.5) {\scriptsize \y};
  \fill[blue!60!black] (5.5,2.5) circle (2pt);
  \node[above right] at (5.5,2.5) {$P(x,y)$};
  \node at (4,-1.45) {Tọa độ hóa từng ô của ma trận.};
\end{tikzpicture}
\end{center}

Bây giờ xét điều kiện để ô $P(x,y)$ có lỗ. Vì ma trận được định nghĩa đệ quy
bằng cách chia đôi, ta tự nhiên chuyển $x$ và $y$ sang nhị phân. Giả sử biểu
diễn nhị phân đủ $N$ bit của chúng lần lượt là
\[
  x=(a_1a_2\cdots a_N)_2,\qquad
  y=(b_1b_2\cdots b_N)_2.
\]
Xét bit đầu tiên $a_1,b_1$. Có bốn cách lấy giá trị, tương ứng với bốn phần
tư trên trái, dưới trái, trên phải, dưới phải trong định nghĩa đệ quy. Khi
$a_1=0$ và $b_1=1$, bất kể các bit phía sau ra sao, ô $P$ không thể có lỗ vì
nó đã rơi vào phần tư trên trái không có lỗ. Trong các trường hợp còn lại,
ta tiếp tục xét bit thứ hai, rồi bit thứ ba, và cứ thế.

\begin{center}
\begin{tikzpicture}[x=1.15cm,y=1.15cm, every node/.style={font=\small}]
  \draw[thick] (0,0) rectangle (4,4);
  \draw[thick] (2,0) -- (2,4);
  \draw[thick] (0,2) -- (4,2);
  \fill[gray!20] (0,2) rectangle (2,4);
  \node at (1,3) {$a_1=0,\ b_1=1$};
  \node at (3,3) {$a_1=1,\ b_1=1$};
  \node at (1,1) {$a_1=0,\ b_1=0$};
  \node at (3,1) {$a_1=1,\ b_1=0$};
  \node at (2,-0.45) {Phân bố bốn trường hợp của $(a_1,b_1)$.};
\end{tikzpicture}
\end{center}

Kết luận là: ô $P(x,y)$ có lỗ khi và chỉ khi không tồn tại vị trí
$1\le i\le N$ sao cho
\[
  a_i=0\quad\text{và}\quad b_i=1.
\]
Ta gọi đây là tính chất có lỗ của một ô. Nói bằng phép toán bit, các bit bằng
$1$ của $y$ phải nằm trong các vị trí mà $x$ cũng bằng $1$.

\paragraph{Giải bằng quy hoạch động.}
Phần sau trở nên rất đơn giản. Điều đề bài cần tìm chỉ là số cặp có thứ tự
$(x,y)$ thỏa
\[
  0\le x,y,\qquad x+X\le 2^N-1,\qquad y+Y\le 2^N-1,
\]
và cả hai ô $(x,y)$, $(x+X,y+Y)$ đều có tính chất có lỗ. Gọi đó là tính chất
có lỗ chung.

Chuyển $X,Y$ sang dạng nhị phân đủ $N$ bit:
\[
  X=(p_1p_2\cdots p_N)_2,\qquad
  Y=(q_1q_2\cdots q_N)_2.
\]
Ta dùng quy hoạch động theo từng bit, ghi nhớ các bit nhớ của hai phép cộng
để bảo đảm trạng thái không phụ thuộc quá khứ xa hơn. Xử lý từ bit thấp lên
bit cao. Đặt
\[
  D_i(c_x,c_y)
\]
là số cách chọn các bit từ vị trí $i$ đến $N$ sao cho mọi vị trí đã xử lý
đều thỏa tính chất có lỗ chung, và sau khi cộng với các bit tương ứng của
$X,Y$, bit nhớ truyền sang vị trí $i-1$ lần lượt là $c_x,c_y$. Mỗi bit nhớ
chỉ có thể là $0$ hoặc $1$.

Khởi tạo $D_{N+1}(0,0)=1$. Với $i=N,N-1,\ldots,1$, từ một trạng thái có bit
nhớ vào vị trí $i$ là $(c_x,c_y)$, thử bốn cách chọn $(a_i,b_i)\in\{0,1\}^2$.
Tính
\[
\begin{aligned}
  u&=a_i+p_i+c_x,& a_i'&=u\bmod 2,& c_x'&=\lfloor u/2\rfloor,\\
  v&=b_i+q_i+c_y,& b_i'&=v\bmod 2,& c_y'&=\lfloor v/2\rfloor.
\end{aligned}
\]
Lựa chọn này hợp lệ nếu đồng thời
\[
  \neg(a_i=0\ \text{và}\ b_i=1)
  \quad\text{và}\quad
  \neg(a_i'=0\ \text{và}\ b_i'=1).
\]
Khi đó cộng số cách vào trạng thái mới mang bit nhớ $(c_x',c_y')$. Sau khi
xử lý hết $N$ bit, đáp án là trạng thái có bit nhớ tràn bằng $0$ ở cả hai
tọa độ; điều này đồng thời bảo đảm $x+X$ và $y+Y$ vẫn nằm trong ma trận.

Mỗi bước chỉ thử $4$ cặp bit cho $4$ trạng thái nhớ, nên nếu bỏ qua chi phí
số nguyên lớn, thời gian là $O(N)$. Dùng mảng luân phiên, không gian chỉ cần
một số hằng định các số nguyên lớn. Chương trình cài đặt cũng rất ngắn. So
với hướng phân tích trực tiếp hình vẽ, ưu thế của ``số'' là rất rõ.

\paragraph{Tiểu kết.}
Khi nhận thấy hình giao của hai ma trận phức tạp và không thuận lợi để xét
trường hợp, ta tránh va chạm trực diện với hình vẽ và chuyển sang mặt số.
Ngay sau đó thu được hai khái niệm rõ ràng: tính chất có lỗ và tính chất có
lỗ chung. Yêu cầu của đề bài vì thế trở nên đơn giản hơn nhiều, và phần còn
lại là một quy hoạch động kinh điển.

Có thể nói đây là một ví dụ tốt, dù những ví dụ thuần túy về dùng số hỗ trợ
hình có vẻ không nhiều. Thực ra, như đã nói, máy tính thông thường đều lấy
số làm nền tảng; khi viết mọi loại chương trình, cuối cùng ta vẫn phải quy
về ``số'' để cài đặt, nên đôi khi vai trò của số trở nên quá quen thuộc đến
mức bị xem nhẹ.

Ví dụ trên cũng cho thấy nếu đề bài có sự định hướng thị giác đủ mạnh, thí
sinh vẫn có thể rời xa con đường ngắn dựa trên số và đi vòng rất xa. Vì vậy,
khi gặp những bài giống ví dụ này, đối diện một hình vẽ phức tạp và đa dạng,
biến hình về số thường là cách tốt để nắm phần then chốt của đề.

\section{Kết luận}

Số và hình là sự trừu tượng hóa và phản ánh các sự vật khách quan trong thế
giới thực. Chúng là nền tảng của toán học, đồng thời cũng là hai phương diện
chính mà đề thi tin học thường chạm tới. Kết hợp số và hình là một tư tưởng
toán học cổ xưa; Olympic Tin học non trẻ lại đem đến cho nó sức sống mới.

Ba ví dụ ở trên, nhìn tổng thể, đều vận dụng khéo léo tư tưởng kết hợp số và
hình. Nhưng xét kỹ hơn, giữa chúng vẫn có vài khác biệt.

Thứ nhất, ba ví dụ trình bày nội hàm của tư tưởng này theo hai điểm nhấn khác
nhau: dùng hình hỗ trợ số và dùng số hỗ trợ hình. Trong bài toán thực tế, số
và hình không có ranh giới tuyệt đối; tư tưởng kết hợp số và hình cũng không
chỉ giới hạn trong hai mặt vừa nêu. Thường gặp hơn là số và hình thúc đẩy
nhau, chứa nhau, và trong điều kiện nhất định chuyển hóa cho nhau. Có thể gọi
đó là ``số hình tương trợ''. Điều này thể hiện quan hệ mâu thuẫn biện chứng
giữa số và hình, cũng như tính đa dạng của tư tưởng kết hợp số và hình.

Thứ hai, dùng ``hình'' để giải ví dụ thứ hai dường như là lối ra duy nhất.
Nhưng với ví dụ thứ nhất và thứ ba, ngoài phương pháp trong bài còn có những
cách giải hay khác. Dù vậy, so sánh tương đối cho thấy việc vận dụng khéo léo
tư tưởng kết hợp số và hình có thể hạ thấp đáng kể độ phức tạp tư duy và lập
trình, mở cho ta một con đường thuận lợi để giải bài trong thời gian thi đấu
ngắn ngủi.

Cũng cần chỉ ra rằng mỗi người có cấu trúc tri thức, kinh nghiệm thi đấu và
cảm nhận về độ khó thuật toán khác nhau. Vì vậy với cùng một bài, những người
khác nhau có thể chọn những con đường số-hình khác nhau. Đây là biểu hiện
của khác biệt cá nhân trong tư tưởng kết hợp số và hình.

Ba ví dụ trong bài đều chọn những thuật toán mà đa số người học có thể tiếp
nhận, nhưng không thể nói đó chắc chắn là thuật toán đơn giản nhất trong cảm
nhận của mọi độc giả. Mục đích của bài viết không nằm ở vài ví dụ nhỏ, mà ở
việc bàn cách vận dụng tư tưởng kết hợp số và hình trong thi tin học.

Vận dụng tư tưởng kết hợp số và hình trong thi tin học nghĩa là khi xử lý bài
toán, ta cân nhắc tình huống cụ thể, nắm mâu thuẫn chính của vấn đề: với bài
toán nặng về quan hệ số lượng, dùng hình học làm nó trực quan và có hình
tượng; với bài toán nặng về hình vẽ, dùng quan hệ số lượng để nắm bản chất.

Kết hợp số và hình, trọng tâm nằm ở chữ ``kết hợp''. Muốn vận dụng linh hoạt
tư tưởng này, cần coi trọng khác biệt cá nhân trong tư duy. Tùy theo trình độ
tri thức và cách suy nghĩ hiện có của mỗi người, hãy kết hợp một cách hữu cơ
toán học trừu tượng, ngôn ngữ máy tính và hình vẽ trực quan; kết hợp tư duy
trừu tượng với tư duy hình tượng; tạo được liên hệ và chuyển hóa giữa khái
niệm trừu tượng và hình ảnh cụ thể. Khi làm được như vậy, ta có thể giải
quyết bài toán thực tế nhanh hơn, tốt hơn và đơn giản hơn.

\appendix
\section{Bài \textit{Sequence} của kỳ tuyển chọn Thượng Hải năm 2003}

\subsection*{Phát biểu bài toán}

Một dãy
\[
  \{A_i\mid i=0,1,2,\ldots,3N\}
\]
có $3N+1$ phần tử, mỗi phần tử bằng $1$ hoặc $-2$. Định nghĩa tổng tiền tố
\[
  S_K=A_0+A_1+\cdots+A_K.
\]
Hỏi có bao nhiêu dãy $A$ thỏa tính chất $P$ sau:
\[
  S_{3N}=1
  \quad\text{và}\quad
  S_K>0\ \text{với mọi }K=0,1,\ldots,3N.
\]
Nói cách khác, tổng mọi phần tử bằng $1$ và mọi tổng tiền tố đều dương.

Ví dụ với $N=2$, có ba dãy:
\[
\begin{array}{l}
1,1,1,-2,1,1,-2,\\
1,1,1,1,-2,1,-2,\\
1,1,1,1,1,-2,-2.
\end{array}
\]
Giới hạn: $N\le 1000$.

\subsection*{Phân tích}

\paragraph{Bổ đề.}
Với mọi dãy $A$ thỏa tổng bằng $1$, không có biểu diễn vòng không tầm thường
nào của $A$ trùng với chính nó.

\paragraph{Chứng minh.}
Nếu có một phép quay không tầm thường tạo lại đúng dãy ban đầu, theo tính
chất của xâu vòng, dãy ấy phải tách được thành $d>1$ phần hoàn toàn giống
nhau. Gọi tổng mỗi phần là $s$. Khi đó $sd=1$. Vì $d>1$, $s$ không thể là số
nguyên, mâu thuẫn với việc mọi phần tử của dãy là số nguyên. Bổ đề được chứng
minh.

\paragraph{Định lý.}
Số dãy thỏa tính chất $P$ là
\[
  \frac{1}{3N+1}\binom{3N+1}{N}.
\]

\paragraph{Chứng minh.}
Trước hết liệt kê mọi dãy có $3N+1$ phần tử, mỗi phần tử bằng $1$ hoặc $-2$
và tổng bằng $1$. Nếu có đúng $m$ phần tử bằng $-2$, tổng của dãy là
\[
  (3N+1-m)-2m=3N+1-3m.
\]
Điều kiện tổng bằng $1$ cho ta $m=N$. Vì thế số dãy như vậy là
\[
  \binom{3N+1}{N}.
\]

Phân loại các dãy theo các biểu diễn vòng của chúng. Nhờ bổ đề, mỗi lớp có
đúng $3N+1$ dãy. Do đó có tất cả
\[
  \frac{1}{3N+1}\binom{3N+1}{N}
\]
lớp. Theo bổ đề Raney, trong mỗi lớp có đúng một dãy có mọi tổng tiền tố
dương. Vậy tổng số dãy thỏa tính chất $P$ chính là
\[
  \frac{1}{3N+1}\binom{3N+1}{N}.
\]

Cùng phương pháp này cũng suy ra công thức của số Catalan.

\section{Ghi chú về phụ kiện của bài gốc}

Bài gốc nhắc tới một báo cáo lời giải cho \textit{Painter's Studio} và một
chương trình Pascal đi kèm. Bản dịch này đã chuyển phần lập luận cần thiết
thành mô tả thuật toán và công thức trong chính bài viết; các tệp phụ kiện
không phải một phần của bản LaTeX hiện tại.

\section{Tài liệu tham khảo}

\begin{enumerate}
  \item Ronald L. Graham, Donald E. Knuth, Oren Patashnik,
  \textit{Concrete Mathematics}.
  \item USA Computing Olympiad, \url{http://ace.delos.com/usacogate}.
  \item Polish Olympiad in Informatics, \url{http://www.oi.edu.pl/}.
  \item Đề thi kỳ tuyển chọn NOI Thượng Hải năm 2003.
  \item Một bài viết về cách dùng tư tưởng kết hợp số và hình trong dạy học
  toán.
\end{enumerate}

\end{document}
